% Split from 3_sylvia.tex by cut-sheet v2/v3 — content below is the source scene, header adjusted
\chapter{Finding the hyperbolic amidst Primitive Triples}
\vspace*{-.8cm}
% \centerline{\textbf{The Geometry of Triples}}
% \section*{}
\begin{minipage}[t]{\textwidth}

\lettrine{S}{ylvia} Kovalevskaya sat before her window, her muse Hector mumbling interminably behind her. His omnipresence was beginning to wear. 

\begin{figure}[H]
\includegraphics[width=\linewidth]{Illustrations/vign-sophie.png}
\caption{Sylvia appealing to the cameraman to kill Hector}
\end{figure}
Just because his dad had fallen fowl of bird flu it didn't mean she was beholden to his every waking breath. In and out. Day after day after day.
\end{minipage}

Consumed by bitterness and regret she could only now ponder prosaic matters such as Pythagorean triples. Those ancient, primitive solutions to 
\[
q^2 + r^2 = s^2.
\]

\begin{figure}[H]
\includegraphics[width=\linewidth]{Illustrations/circleHomo.png}
\caption{Intersected circle with its significant homogeneous co-ordinates.}
\end{figure}

And then... her hand hovered over her keyboard as she began to think of these triples not just as static points, but as dynamic entities born of geometric interactions. Sylvia instead reached for her pad and wrote out the equation for the unit circle in cartesian and polar form:
\[
x^2 + y^2 = 1,\quad cos^2\theta+sin^2\theta=1.
\]
Trigonometric functions described the periodicity, cycles of sine and cosine tracing the circle’s edge.
Consider now the pencil of lines whose gradients are their y-intercepts:
\(
y = mx + m.
\) which when combined, have intersection points that are Pythagorean triples:
\[
x^2 + m^2(x + 1)^2 = 1.
\]
Here, the line isn’t just a line it is a trajectory slicing through the periodicity of the circle. 

\begin{figure}[H]
\includegraphics[width=\linewidth]{Illustrations/logAreaAcutePPTs_1000labelled.png}
\caption{The bead curtain of 30000 triples with detractor primitives highlighted.}
\end{figure}

The intersection points represented a moment of interaction, the crossing of two infinitesimals in space. \emph{Two particle paths,} she thought, \emph{two trajectories weaving through spacetime.} Just a bullet, tracing a parabolic path cutting through tissue,...no not now.

Hector, muttered. "Yes my dear that is so".

Sylvia flipped to a page in her notebook and sketched a hyperbola, the curve defined by:
\[
x^2 - y^2 = 1.
\]
It was the free spirited sibling to the closed circle, but stretched, opened: a mirror of infinite trajectories. She jotted:
\begin{quote}
The circle, bound and closed, represents periodicity: cycles, oscillations, the sine and cosine. The hyperbola, unbound, of sinh and cosh of growth and decay.
\end{quote}

% >>>>>> ANNEX B (cut-sheet v2): Pythagorean triples re-parameterised into hyperbolic geometry — block commented out below, pointer left in text <<<<<<
\textit{(The full working --- circle to hyperbola, triples to rapidities --- is reproduced in Annexe B; what mattered tonight was the shape of the thought.)}

% \section*{From Pythagorean Triples to Hyperbolic Geometry}
% 
% Sylvia pondered aloud, "From the Pythagorean theorem on the Euclidean plane, why not go hyperbolic?" and consider hyperbolic triples \(q^2-r^2=s^2\) satisfying:
% \[
% x^2 - y^2 = 1 = (x-y)(x+y).
% \]
% "Ah,but these are \textit{just} the difference of two squares." And such a square unit hyperbola can be parameterized by the hyperbolic identity:
% \begin{equation*}
% \cosh^2\phi - \sinh^2\phi = 1. \label{hyperbolic}
% \end{equation*}
% "For straight lines \(y = mx + b\), with slope \(m = \tan(\phi)\), intersecting the hyperbola we have:
% \[
% x^2 - (mx + b)^2 = 1.
% \]
% Expanding and rearranging yields:
% \[
% (1 - m^2)x^2 - 2mbx - (b^2 + 1) = 0.
% \]
% For the special case \(b = 0\), the intersections reduce to:
% \[
% x = \pm\sqrt{\frac{1}{1 - m^2}}, \quad y = m \cdot \sqrt{\frac{1}{1 - m^2}}.
% \]
% "These," Sylvia uttered to Hector, "are the 'triple' signatures of hyperbolic space. And these rather remind me when \(m = v/c\) of my Ernest and his fondness for the Lorentz gamma factor, where \(v\) is the relative speed as a fraction of the speed of light." "Now," she added, "if we consider the case where \(b = m\). Substituting \(b = m\) into the equation gives:
% \(
% (1 - m^2)x^2 - 2m^2x - (m^2 + 1) = 0.
% \)
% which we solve Hector as you even well know using the quadratic formula:
% \(
% x = \frac{-B \pm \sqrt{B^2 - 4AC}}{2A},
% \)
% where \( A = (1 - m^2) \), \( B = -2m^2 \), and \( C = -(m^2 + 1) \).
% \[
% x = \frac{2m^2 \pm \sqrt{4}}{2(1 - m^2)} = \frac{2m^2 \pm 2}{2(1 - m^2)}.
% \]
% The solutions for \(x\) are:
% \(
% x_1 = \frac{m^2 + 1}{1 - m^2}, \quad x_2 = -1.
% \)
% For \(y = mx + m\), substitute \(x_1\) and \(x_2\):
% \[
% y_1 = m \cdot \frac{m^2 + 1}{1 - m^2} + m = \frac{m(m^2 + 1) + m(1 - m^2)}{1 - m^2} = \frac{2m}{1 - m^2},
% \]
% \[
% y_2 = m \cdot (-1) + m = 0.
% \]
% Thus, the solutions are:
% \(
% (\frac{m^2 + 1}{1 - m^2}, \frac{2m}{1 - m^2});(-1,0).
% \)
% Considering the case where \(m = \frac{1}{2}\) we have:
% \[
% x_1 = \frac{\left(\frac{1}{2}\right)^2 + 1}{1 - \left(\frac{1}{2}\right)^2} = \frac{\frac{1}{4} + 1}{1 - \frac{1}{4}} = \frac{\frac{5}{4}}{\frac{3}{4}} = \frac{5}{3}, \quad
% y_1 = \frac{2 \cdot \frac{1}{2}}{1 - \left(\frac{1}{2}\right)^2} = \frac{1}{\frac{3}{4}} = \frac{4}{3}.
% \]
% Thus, for \(m = \frac{1}{2}\), the solutions are ( written in homogeneous coordinates of projective geometry) as:
% \(
% [x_1 : y_1 : 1] = [5 : 4 : 3].
% \)
% \begin{center}
% \begin{figure}[H]
% \includegraphics[width=0.8\linewidth]{Illustrations/unitToHyperbolaMAp.png}
% \caption{Pythagorean Triples from Unit Circle and Hyperbola}
% \end{figure}
% \end{center}
% 
% \subsection*{Hyperbolic space-time}
% 
% The hyperbolic nature of space-time transformations in special relativity hinges on kindred principles. The gamma factor is present in the inertial rest mass-energy equation of Einstein \(E=mc^2=\gamma m_0c^2\) is given by:
% \begin{equation*} 
% \gamma(v) = \frac{1}{\sqrt{1 - v^2/c^2}}\label{gamma}.
% \end{equation*}
% Squaring this we can equivalently write it as hyperbolic trig identity, (~\ref{hyperbolic}):
% \[
% \gamma^2 - \left(\gamma \frac{v}{c}\right)^2=1= \cosh^2\phi - \sinh^2\phi,
% \]
% by relating hyperbolic functions with \textit{rapidity} \(\phi\) to boosts according to
% \[
% \cosh\phi = \gamma \quad \text{and} \quad \sinh\phi = \gamma \frac{v}{c}.
% \]
% 
% \noindent
% Now consider again the straight line \(y = mx + b\), for \(b=0\) in which the gradient is parametrized in terms of \(m = v/c\) and whose intersections with the hyperbola \(x^2 - y^2 = 1\) yield:
% \[
% x = \pm\frac{1}{\sqrt{1 - v^2/c^2}}= \pm \cosh\phi, \quad y = \frac{v}{c}\cdot \frac{1}{\sqrt{1 - v^2/c^2}}\sinh\phi.
% \]
% \noindent
% We have that the relative velocity \(v\) is:
% \(
% v = c \tanh\phi\) , for \(\tanh\phi = \frac{\sinh\phi}{\cosh\phi}
% \)
% \begin{figure}[H]
% \includegraphics[width=0.7\linewidth]{Illustrations/tanh.png}
% \caption{tanh function}
% \end{figure}
% 
% Sylvia still abusing her muse, mutters "Both hyperbolic geometry and Lorentz transformations are rooted in the invariance of hyperbolic distances." She paused, the ideas a lingering in the air. "So perhaps these hyperbolic triples are but reflections of a deeper, universal symmetry."
% 
% She wrote: \emph{Just as we find Pythagorean triples in the flat plane, so too can we find their analogs in the saddle of hyperbolic geometry. But how do these play out in Minkowski or Lobachevsky spaces with negative signature?}
% 
% Sylvia paused, imagining a Minkowski space proper now: a geometry where time and space entwined. The familiar Pythagorean equation
% \(
% x^2 + y^2 = z^2,
% \) morphed relativistically into:
% \[
% x^2 + y^2 - t^2 = s^2,
% \]
% where \(t\) was time and \(s\) the spacetime interval. 
% 
% \emph{Here, triples would evolve,} she imagined. \emph{Instead of merely measuring Euclidian lengths on the flat, they describe connections between events in spacetime}; the triples as quadruples would become threads of causality, vertices of a spacetime lattice of triangles.
% 
% \begin{figure}[H]
% \includegraphics[width=\linewidth]{Illustrations/vign-Sophia2.png}
% \caption{Sylvia getting ahead of herself}
% \end{figure}
% 
% But what of the Lobachevsky geometry? The saddle space of a curved plane warped like a horse’s saddl, offered new challenges in which the interplay of geodesics, the straightest possible paths on a curved surface, created analogs to Pythagorean triples. 
% 
% Sylvia began to sketch. She imagined a geodesic triangle on the saddle, its vertices connected by paths that weren’t straight in the Euclidean sense but curved, bowing inward under the pull of negative curvature. The lengths of these sides could still satisfy an analog to Pythagoras, but where hyperbolic trigonometric functions described the measures. \emph{What of integer solutions here?} she wondered. \emph{Can there exist ‘hyperbolic triples’ in the Lobachevsky geometry, bound by such warped relationships?} 
% 
% \emph{"Can hyperbolic tuples be the key to understanding wave behavior in curved space?"} she murmured. Just as light refracted through a prism, disperses into its spectrum, so too might hyperbolic waves on some lattice, some triangluar lattice, reveal spectral features as the saddle’s curvature does the job of Snell. 
% 
% But perhaps, she thought, the answer isn’t in finding the measures of these triples, but as is often the case in maths don't dwell on the triples rather explore the paths and trajectories that connect them.
% 
% \begin{table}[H]
% \centering
% \renewcommand{\arraystretch}{1.5}
% \setlength{\tabcolsep}{8pt}
% \begin{tabular}{|c|c|c|c|}
% \hline
% \textbf{\(m\)} & \textbf{Intersection Co-ord} & \textbf{\(\theta\)} & \textbf{Pythagorean Triple} \\
% \hline
% \(1/11\) & \((61/60, 11/60)\) & 0.18 & \(60^2 + 11^2 = 61^2\) \\
% \(1/10\) & \((101/99, 20/99)\) & 0.2 & \(99^2 + 20^2 = 101^2\) \\
% \(1/9\) & \((41/40, 9/40)\) & 0.22 & \(40^2 + 9^2 = 41^2\) \\
% \(1/8\) & \((65/63, 16/63)\) & 0.25 & \(63^2 + 16^2 = 65^2\) \\
% \(1/7\) & \((25/24, 7/24)\) & 0.29 & \(24^2 + 7^2 = 25^2\) \\
% \(1/6\) & \((37/35, 12/35)\) & 0.34 & \(35^2 + 12^2 = 37^2\) \\
% \(2/11\) & \((125/117, 44/117)\) & 0.37 & \(117^2 + 44^2 = 125^2\) \\
% \(1/5\) & \((13/12, 5/12)\) & 0.41 & \(12^2 + 5^2 = 13^2\) \\
% \(2/9\) & \((85/77, 36/77)\) & 0.45 & \(77^2 + 36^2 = 85^2\) \\
% \(1/4\) & \((17/15, 8/15)\) & 0.51 & \(15^2 + 8^2 = 17^2\) \\
% \(3/11\) & \((65/56, 33/56)\) & 0.56 & \(56^2 + 33^2 = 65^2\) \\
% \(2/7\) & \((53/45, 28/45)\) & 0.59 & \(45^2 + 28^2 = 53^2\) \\
% \(3/10\) & \((109/91, 60/91)\) & 0.62 & \(91^2 + 60^2 = 109^2\) \\
% \(1/3\) & \((5/4, 3/4)\) & 0.69 & \(4^2 + 3^2 = 5^2\) \\
% \(4/11\) & \((137/105, 88/105)\) & 0.76 & \(105^2 + 88^2 = 137^2\) \\
% \(3/8\) & \((73/55, 48/55)\) & 0.79 & \(55^2 + 48^2 = 73^2\) \\
% \(2/5\) & \((29/21, 20/21)\) & 0.85 & \(21^2 + 20^2 = 29^2\) \\
% \(3/7\) & \((29/20, 21/20)\) & 0.92 & \(20^2 + 21^2 = 29^2\) \\
% \(4/9\) & \((97/65, 72/65)\) & 0.96 & \(65^2 + 72^2 = 97^2\) \\
% \(5/11\) & \((73/48, 55/48)\) & 0.98 & \(48^2 + 55^2 = 73^2\) \\
% \(1/2\) & \((5/3, 4/3)\) & 1.1 & \(3^2 + 4^2 = 5^2\) \\
% \hline
% \end{tabular}
% \caption{Intersection Coordinates and Pythagorean Triples for Various Gradients}
% \label{tab:pythagorean_triples}
% \end{table}
% 
% \section*{postscript}
% 
% \begin{figure}[H]
% \includegraphics[width=\linewidth]{Illustrations/pythagHyperbola.png}
% \caption{There is not unique map from gradients to triples}
% \end{figure}
% \sloppybottom
% 
% >>>>>> end of annexed block <<<<<<
% >>>>>> ANNEX E (cut-sheet v2): from Pythagorean triples to elliptic curves (the ECC on-ramp) — block commented out below, pointer left in text <<<<<<
\textit{(Her route from triples to elliptic curves --- the same curves that would one day underwrite public-key cryptography --- is set out in Annexe E.)}

% \section*{Pythagorean Elliptic Curves}
% 
% A Pythagorean triple \((a, b, c)\) satisfies:
% \(
% a^2 + b^2 = c^2.
% \)
% From this, we define the normalized rational coordinates:
% \[
% u = \frac{a}{c}, \quad v = \frac{b}{c}.
% \]
% For each triple, we seek an \textbf{elliptic curve} of the form:
% \(
% v^2 = u^3 + k u.
% \)
% Solving for \( k \) using the given \((u,v)\), we obtain a different \( k \) for each triple. This code \href{https://colab.research.google.com/drive/1_MGQT3AjpWPs7RNE2Q_kPcKp3P7Lfgtx?usp=sharing}{ellipticPythag.ipynb} gives us the table of the first 20 Pythagorean triples and their associated elliptic curves.
% 
% \begin{table}[h]
%     \centering
%     \begin{tabular}{ccccc}
%         \toprule
%         $a$ & $b$ & $c$ & $d^2$ & Elliptic Coefficient ($k$) \\
%         \midrule
%         3  & 4  & 5  & 36     & $\frac{53}{75}$ \\
%         5  & 12 & 13 & 900    & $\frac{1747}{845}$ \\
%         7  & 24 & 25 & 7056   & $\frac{14057}{4375}$ \\
%         % 8  & 15 & 17 & 3600   & $\frac{3313}{2312}$ \\
%         % 9  & 40 & 41 & 32400  & $\frac{64871}{15129}$ \\
%         \bottomrule
%     \end{tabular}
%     \caption{Pythagorean Triples and Corresponding Elliptic Coefficients}
%     \label{tab:pythagorean_elliptic}
% \end{table}
% For the \((3,4,5)\) triple:
% \(
% u = \frac{3}{5}, \quad v = \frac{4}{5},
% \)
% we find the elliptic curve:
% \begin{equation}  
% v^2 = u^3 + \frac{53}{75} u. \label{elliptic345}
% \end{equation}
% 
% To generate additional rational points, we apply the tangent (to \ref{elliptic345} so that \(2\frac{dv}{dt}=3u^2+\frac{53}{70}\)) method where the slope at \((u_0, v_0)\) is given by:
% \begin{equation}  
% m = \frac{3 u_0^2 + \frac{53}{75}}{2 v_0}. \label{2elliptic345}
% \end{equation}
% The equation of the tangent line is:
% \(
% v = m (u - u_0) + v_0.
% \)
% Solving for the intersection of this line with the elliptic curve yields a new rational point \((u_1, v_1)=(\frac{43}{916}, \frac{110}{603}\)).
% 
% \begin{figure}[H]
% \includegraphics[width=\linewidth]{Illustrations/ellipticTangent.png}
% \caption{Tangent metod to generate new triples}
% \end{figure}
% 
% To iteratively generate additional rational points \(.(u_1, v_1), (u_2, v_2), (u_3, v_3),..\), rather than relying on solving for intersections directly, we use the \textbf{group law} for elliptic curves .  Given two points \( P = (u_1, v_1) \) and \( Q = (u_2, v_2) \): If \( P = Q \), the \textbf{point doubling formula} is applied \ref{2elliptic345}, otherwise, for distinct points, the \textbf{addition formula} is used:
%   \(
%   m = \frac{v_2 - v_1}{u_2 - u_1}.
%   \). The new point is given by:
%   \(
%   u_3 = m^2 - u_1 - u_2, \quad v_3 = m (u_1 - u_3) - v_1.
%   \)
% 
% \begin{table}[h]
%     \centering
%     \begin{tabular}{cc}
%         \toprule
%         $u$ & $v$ \\
%         \midrule
%         $\frac{3}{5}$     & $\frac{4}{5}$ \\
%         $\frac{43}{916}$  & $-\frac{110}{603}$ \\
%         $\frac{148}{59}$  & $-\frac{1609}{384}$ \\
%         % $\frac{1070}{287}$ & $\frac{6088}{825}$ \\
%         % $\frac{95}{996}$  & $\frac{173}{662}$ \\
%         \bottomrule
%     \end{tabular}
%     \caption{Generated Rational Points on the 3-4-5 Elliptic Curve}
%     \label{tab:elliptic_rational_points}
% \end{table}
% 
% \begin{tcolorbox}[title=Elliptic Curve Point Addition, colback=blue!5!white, colframe=blue!75!black]
% Given two distinct rational points \( P = (u_1, v_1) \) and \( Q = (u_2, v_2) \) on an elliptic curve of the form \( v^2 = u^3 + ku \), the \textbf{addition formula} yields a third point \( R = (u_3, v_3) \), defined by:
% 
% \[
% m = \frac{v_2 - v_1}{u_2 - u_1}, \quad
% u_3 = m^2 - u_1 - u_2, \quad
% v_3 = m (u_1 - u_3) - v_1.
% \]
% 
% This rule arises from intersecting the elliptic curve with the line through \( P \) and \( Q \), and taking the reflection of the third intersection point across the horizontal axis to preserve group closure.
% 
% If \( P = Q \), use the \textbf{point doubling formula} instead:
% 
% \[
% m = \frac{3u_1^2 + k}{2v_1}, \quad
% u_3 = m^2 - 2u_1, \quad
% v_3 = m (u_1 - u_3) - v_1.
% \]
% \end{tcolorbox}
% 
% >>>>>> end of annexed block <<<<<<
\section*{Idempotents}

The café was almost empty. Ernest arrives early, as he always did,
and had already aligned the sugar packets with the edge of the table
by the time Sylvia entered. They sat opposite one another in the winter light of the café. Ernest caught himself stirring his coffee with unnecessary thoroughness. She always made him nervous

“You’re still working on quadratic structures?” he asked,
as though they had spoken yesterday.

“Always,” she said. “They don’t exhaust themselves.”

“You know,” he began, “idempotents are reassuring.
Apply them twice and nothing changes.
Stability is built into their algebra.”

Sylvia watching his spoon circle.
“Or paralysis,” she said.
“Nothing changes because nothing can.”

Smiling faintly. “That’s unfair.
Idempotents carve space cleanly, so splitting structure into components.”

“When it splits,” she replied.
“And when it doesn’t?”

Getting into it. “Then you don’t get projectors.
You get coupling.”

“You get inertia,” she retorts.
“A field that won’t decompose.”

One hiatus follows another.
Outside, a bus exhaling at the curb.

“Nilpotents are worse,” Ernest offers darkly.
“Power them up and they disappear.”

Sylvia’s fingers tracing the rim of her cup impotently.
“Nilpotents are honest.
They admit collapse.
They don’t pretend to persist.”

Ernest’s expression didn’t change.
“Most systems aren’t nilpotent.”

“It’s structural. A nilpotent tells you in advance
that iteration will not produce life.”

“That’s a bleak way to put it.”

“It’s accurate,” 
“Some elements square to zero.
No matter how carefully you structure them.”

Ernest shifting - always aware of an undercurrent.
“You always preferred degeneracy theory.”

“I prefer knowing when something is irreversibly altered.”

Looking at her then, good and properly.
“Involution is cleaner,” 
“Apply it twice and you return to where you started.”

Sylvia now holding his glaze.
“Involution assumes you can go back.”

Her words rested on his chest.

Ernest clearing his throat goes for what he thinks is a dsitractor.
“Ramification is messier.
An ideal that doesn’t split.
Multiplicity where you expected two points.”

“Yes,” says Sylvia.
“A single point thickened.
The square of something already fragile.”

Ernest now ernestly nodding. “But ramification is controlled.
It’s finite. It doesn’t spread.”

“It spreads enough,” she answers soulfully.
“You live with its consequences in every fibre.”

He knew now was not the time to respond.

But such is his irrepressibility he tries again.
“The global invariant doesn’t care about local accidents.
It integrates everything. Balances the signs.”

“Does it?” she asks. "Does it really, Ernest-does it?"
“Or does it record what fails to cancel?”

Ernest floundering.
“That’s the index. Yes the index that we are talking about ... isn't it?”

“Yes,” Sylvia releases him.
“The excess.”

They sit with that.

Outside, the winter light is thinning towards evening. Inside, nothing is moving except the cooling particulates of escaping coffee.

“Some structures return,” Ernest says at last.
“Some remain.”

“And some vanish,” Sylvia replied.
“What doesn’t cancel is what you carry.”

\section*{Ramifications}

He nods. “Idempotents, then.
Clean things. Apply them twice and nothing further happens. They stabilise,” Sylvia agreed. But Nilpotents are more troubling. They annihilate themselves. Power them and they vanish.”

“They don’t vanish,” Sylvia replies. “They reveal their order.” while nodding knowingly without knowing and not knowing anyhow. “That’s generous.”

“Most systems don’t admit reversal either.”

Ernest considers that. “Involution does. Apply it twice and you return.”

“Yes,” Sylvia, softly. “Involution presumes symmetry.”

Ernest then reache
The spoon touches porcelain. A small, precise sound.

“Ramification is different,” Ernest said.
“It isn’t collapse and it isn’t symmetry.
An ideal that doesn’t split.
Multiplicity without separation.”

“A thickened point,” Sylvia said.
“One place bearing double weight.”

“It’s local,” he answered evenly.
“It doesn’t contaminate the whole structure.”

“Doesn’t it?”

He met her eyes without flinching.
“In arithmetic, it doesn’t.
You isolate it. Account for it.
The global invariant absorbs the defect.”

Sylvia looked past him, toward the window.
“Absorbs,” she repeated.

“That’s the point of the index.
You integrate local contributions.
What cancels disappears.”

“And what doesn’t?”

He hesitates just a fraction of a second.
“What doesn’t cancel remains in the total.”

“The excess,” she said.

“Yes.”

They sit with a joint affirmative. Outside, traffic moves in disciplined lines. Inside, nothing moves but shifts.

“Not every structure splits cleanly,” Ernest says at last.
“But most of them behave.”

Sylvia smiled, not wholly unkindly.
“Behaviour isn’t the same as innocence.”

He gives a small nod, as though acknowledging their shared loss. She traces a circle in the condensation on her glass.
“Finite doesn’t mean light.”

“No,” he said.
“But it means bounded.”

After a moment he adds, almost academically,
“Ramification is finite.”

The light shifting across the table, leaving one half in shadow.
“Most ideals still factor,” Ernest adds,
as though clarifying a theorem. “Ramification is exceptional.”

“Yes,” Sylvia says, “It always is. And so are consequences.”
\section*{Fibril Times}

Ernest has moved on to speaking about correction factors,
his voice measured, careful.
Sylvia listens, but her mind is elsewhere. Fibres a firing.

Sylvia has been reworking in her mind her Eisenstein picture, of  the elliptic curve sitting above the base, primes rising as points, splitting into two when symmetry permitted, thickening when they would not. She preferred it that way: not as a product over codes, not as encryption and residue classes, but as a map of fibres over a line. She preferred it that way not least because John would have liked this version. Yes, good old  John would have liked it. He had patience for geometry. He would trace the morphism with his finger as though it were terrain, not arithmetic, declaring “Look, there,that’s where the structure lives.” Eleanor would likely have said something else entirely. Eleanor would have asked how it could be encoded, how the splitting law might generate keys, how multiplicative structure hides information. Eleanor liked arithmetic because it sealed things.

Sylvia turned her cup slowly.

In the Eisenstein case, the fibres were clean: two points when $p\equiv1\pmod{3}$,
one when inert, a doubled point when ramified. She likes that the geometry did not lie.
Multiplicity shows itself. That thickening was visible.

Ernest was saying something about functional equations now,
about completion. She catches only the word “absorb.”

John would not say absorb. He would say "carry".

She imagines explaining it to him: that the class number was not a statistic,
but a measure of how many fibres refused to trivialise globally.
That $\Pic(\Spec(\mathcal{O}_K))$ was not abstraction, but rather the memory of obstruction.

Eleanor would call it inefficiency.
John would call it curvature.

Ernest pauses, waiting for an agreement.

“Yes,” Sylvia says aloud, though she was thinking:
A collaboration is controlled but companionship is something else.
It is not encryption. Not involution, not squaring back to a start.

Ramification is finite, she reassures herself. But fibres persist.

Ernest stopping mid-sentence. “You’re somewhere else,” he says quietly.

Sylvia blinks. “I’m here. In the café, yes.”

Flat now. He does not smile when he says it.

Sylivia, intertwining her hands now, knuckles whitening.
“I was thinking about fibres.”

“You prefer that language now.”

“It’s more honest.”

“More honest than what?”

“Than pretending everything reduces to residue classes.”

He studies her face. “That’s not what I do.”

“I know.”

A pause.

“But it’s what Eleanor does,” he said, almost lightly.

Sylvia looks up. “She does it well.”

“And John?”

The name entered the space without emphasis.

Sylvia does not answer immediately. “John likes geometry,” she utters at last "He trusts maps more than ciphers.”

“And you?”

Sylvia, considering the condensation ring her glass now was leaving behind. “In fibre geometry, multiplicity is visible. If something is thickened, you see it. It doesn’t hide behind notation.”

Ernest nods once. “Yes. Geometry exposes structure.”

“Arithmetic encodes it,” she said.

“And you would rather see it?”

“I would rather know where the weight sits.”

Ernest absorbs that without comment.

After a moment reuttering, “Ramification doesn’t disappear in analytic completion. It’s still there. It just contributes differently.”

“Yes,” Sylvia says. “Don't I know it.”

He holds her gaze. “And with John?” 

“With John, I suspect, you would not be discussing encryption.”

Air between them shifts, discernibly, Sylvia  unsmiling now “It wouldn’t be of that encryption stuff.”

“No, No, it wouldn’t.” Silence again.

“Ramification is finite,” Ernest added,
almost reflexively.

Sylvia meeting his eyes again this time.

“Yes,”  “But it changes the fibre.”

Ernest did not look away.

“Fibres can change,” he says evenly. “The base remains. In a morphism,” he replies, “the base is fixed.”

Sylvia exhales through her nose, resting her hands on the table, palms flat, as if steadying something unseen.

“There are constructions,” he says, now more carefully,
“where ramification alters the fibre but the global invariant remains unchanged.”

The winter light has thinned further outside, as the reflections in the glass grow darker, superimposing their faces over the street. Ernest looks down at his hands. Why not? They are reptilian,

\section*{The Hidden Lattice}

The wheel slowed but did not stop entirely. A faint ticking sound came from the ratchet as it settled. Each spoke caught the observatory lamp for a moment and then released it again. The chrome hub reflected Ernest's face in a distorted oval all moustache, breath, and sharp  nose bending around the metal.

“Turn it again,” Ernest says softly.

Sylvia obliges , the wheel moving slowly now. The spokes passed across one another and the pattern emerged once more: not circular this time, but triangular. Then hexagonal. Ernest inhaled sharply. “That’s it.”

Sylvia watched him carefully. “What is?”

He pointed into the wheel. “The \(A_2\) lattice.”

She cocking her head with no expectation to see. “Show me.”

\begin{figure}[H]
\includegraphics[width=\linewidth]{Illustrations/vign-chair.png}
\caption{Sylvia in a better alignment}
\end{figure}

Ernest drew quickly on the back of the observing sheet.

\[
\alpha_1 = (1,0), \qquad 
\alpha_2 = \left(-\tfrac12, \tfrac{\sqrt{3}}{2}\right)
\]
Then he connected them. “These are the simple roots of \(SU(3)\).”
He rotated the paper slightly.
“Sixty degree separation.”

Sylvia nodding slowly, “A triangular lattice.”

“Yes,” as he wrote another expression beneath the roots. “Eisenstein integers.”

\[
a + b\omega, \qquad 
\omega = e^{2\pi i/3}
\]

Sylvia’s eyes follows the symbols. “And the norm.”
Ernest nodded. Together they said it almost simultaneously.
\[
N(a+b\omega) = a^2 - ab + b^2
\]
Sylvia smirking now faintly “Or in the rotated frame…” She taps the equation he had written earlier.
\[
X^2 + XY + Y^2.
\]
Ernest staring at the wheel again. Now that he had seen it he could not un-see it. The spokes divided the circle into equal arcs. Through their crossings the hidden triangular geometry appeared again and again, like a lattice trying to reveal itself through motion. “The wheel is a projection,” he announces unnecessarily slowly but for good effect.

Sylvia raises an eyebrow. “A projection of what?”

Ernest gestured toward the equations. “A hexagonal metric.”

He writes again: \(
g_{\mathrm{eff}} = X^2 + XY + Y^2
\)
“If a hidden sector obeyed this geometry…”

Sylvia finishes the thought quietly. “…then our usual relativistic form would only be an approximation.”

Ernest turns the page back to his earlier dispersion equation.
\(
E^2 = (mc^2)^2 + (pc)^2
\), “The diagonal form.” He looked at Sylvia. “But if the sectors share a lattice structure…”

She nods. “Then a cross term must appear.” Ernest wrote one more line.
\[
E^2 = A^2 + \lambda AB + B^2
\]
The observatory lamp flickers slightly as the wind shifts outside the dome. Sylvia’s wheel moved again beneath her fingers. For a moment the spokes aligned perfectly. The triangular lattice appeared with startling clarity inside the circle. Ernest laughed quietly. 

“I spent ten years staring at galaxy rotation curves, ” shaking his head. “And the geometry was sitting here all along.”

Sylvia watching him. The brittle tension she had carried earlier had softened. “You only see patterns,” she said, “when you stop protecting them.”

Ernest looked puzzled. “What do you mean?”

She gestured toward the equations. “You stopped protecting the old form.” 
\(
E^2 = (mc^2)^2 + (pc)^2
\)
“And suddenly a larger geometry appeared.”

Ernest studied the wheel again. The hexagon dissolved. Only the circle remained. But now the circle looked incomplete. “The dark sector,” he said quietly.

Sylvia nodded. “Another set of coordinates.”

Ernest wrote the final line carefully. Ernest stepping back as Sylvia turns the wheel once more and the spokes cross. The hidden lattice appears. And this time both of them see it clearly.
